\documentclass[10pt,a4paper]{amsart}
\usepackage[margin=19mm]{geometry}
\usepackage{amsmath,amssymb,mathtools}
\usepackage[scr=rsfs,cal=euler]{mathalfa}
\usepackage{xfrac}
\usepackage[unicode,hidelinks,bookmarks=false]{hyperref}
\newcommand{\ie}{i.e.\ }
\newcommand{\cf}{cf.\ }
\newcommand{\resp}{respectively\ }
\newcommand{\Subsection}{\subsection}
\providecommand{\nobreakdash}{\nobreak\hbox{-}}
\setlength{\emergencystretch}{2em}
\multlinegap0pt
\newtheorem{restatable}{Restatable}
\newtheorem{definition}{Definition}
\newtheorem{theorem}{Theorem}
\usepackage{cleveref}
\crefname{restatable}{Restatable}{Restatables}
\crefname{definition}{Definition}{Definitions}
\crefname{theorem}{Theorem}{Theorems}
\newcommand{\R}{\mathbb{R}}
\title{Splitting Sandwiches Unevenly via Unique Sink Orientations and Rainbow Arrangements}
\author{Michaela Borzechowski, Sebastian Haslebacher, Hung P. Hoang, Patrick Schnider, Simon Weber}
\date{Source: arXiv:2602.10795v1 (CC BY 4.0)}
\begin{document}
\maketitle

\noindent\emph{Source and licence.} Splitting Sandwiches Unevenly via Unique Sink Orientations and Rainbow Arrangements. Version arXiv:2602.10795v1 (CC BY 4.0), distributed by arXiv under the Creative Commons Attribution 4.0 International licence, http://creativecommons.org/licenses/by/4.0/, as recorded in the arXiv licence metadata for this exact version and shown on the abstract page https://arxiv.org/abs/2602.10795v1. Full licence notice: \texttt{licenses/arxiv-2602.10795-CC-BY-4.0.txt}.\par\smallskip

\noindent\textit{Editorial scope.} Counted unit: the article's theorem thm:beta-gamma (source0.tex lines 418--421). The article's complete original proof of it is delivered with it (source0.tex lines 423--433, one \texttt{proof} environment of 11 lines). No mathematics is written, repaired or re-derived: the statement and the proof are the article's own lines, in the article's own notation, and no retained line is substituted or re-flowed.  Retained uncounted as the source-local context the proof needs: lines 201--204; lines 404--414.  Nothing was omitted from any retained span, so the package carries no omission record.  No retained line contains a citation, so the wrapper carries no bibliography and no bibliography entry.  No retained line contains a figure environment, an image-inclusion command or drawing code, so no image is delivered and the package declares no asset dependency.  Editorial formatting only, in addition: an AMS \texttt{amsart} preamble that restates the article's own declarations and macros used by the retained lines, namely the environments \texttt{restatable}, \texttt{definition}, \texttt{theorem} on separate counters, together with the standard packages those lines need.  The retained environments print in this document as Restatable 1, Definition 1, Theorem 1 in order of appearance, whereas the article prints the same items with the numbering of its own full document.  The complete native source of the article as distributed in the arXiv e-print for this exact version is bundled unchanged as \texttt{sources/194428/source0.tex}.
\begin{restatable}{theorem}{rainbowHS}\label{thm:rainbowHS}
Let $\mathcal{F}=(H_1, \ldots,H_d)\subset\R^d$ be a well-separated colored generalized arrangement where each color class $H_i$ has size $n_i$. Then for every $(\alpha_1, \dots, \alpha_d) \in [n_1] \times \dots \times [n_d]$, there is a point $x_{\alpha} \in \mathbb{R}^d$ lying on one and above exactly $\alpha_i - 1$ pseudo-hyperplanes of color $i$ for all $i \in [d]$.
If $\mathcal{F}$ is additionally a rainbow arrangement, then the point $x_\alpha$ is unique.
\end{restatable}

\begin{definition}[$(\beta,\gamma)$-separated]
\label{def:generalized_separation}
Let $\mathcal{F}=(H_1, \dots, H_d)$ be a colored generalized arrangement in $\mathbb{R}^d$, where each color class $H_i$ has size $n_i$. Let $\beta=(\beta_1,\ldots,\beta_d)$ and $\gamma=(\gamma_1,\ldots,\gamma_d)$ be vectors such that $\beta_i, \gamma_i \in [n_i]$ and $\beta_i\leq \gamma_i$ for all $i \in [d]$.
The colored generalized arrangement $\mathcal{F}$ is called \emph{$(\beta,\gamma)$-separated} iff for every sign vector $s \in \{+, -\}^d$, there exists a point $x_s \in \mathbb{R}^d$ such that 
\begin{itemize}
	\item $x_s$ lies in an unbounded cell,
    \item $x_s$ lies above the $\gamma_i$-level of $H_i$ iff $s_i = +$,
    \item and $x_s$ lies below the $\beta_i$-level of $H_i$ iff $s_i = -$,
\end{itemize}
for all $i \in [d]$.
\end{definition}

\begin{theorem}
\label{thm:beta-gamma}
Let $\mathcal{F}=(H_1, \ldots,H_d)$ be a $(\beta,\gamma)$-separated colored generalized arrangement where each color class $H_i$ has size $n_i$. Then for every $(\alpha_1, \dots, \alpha_d) \in \{\beta_1,\ldots,\gamma_1\} \times \dots \times \{\beta_d,\ldots,\gamma_d\}$, there is a point $x_{\alpha} \in \mathbb{R}^d$ lying on one and above exactly $\alpha_i - 1$ oriented pseudo-hyperplanes of color $i$ for all $i \in [d]$.
\end{theorem}

\begin{proof}
By assumption $\mathcal{F}$ is $(\beta,\gamma)$-separated, so for each vector $s \in \{+, -\}^d$ there is a point $x_s$ in an unbounded cell that lies above or below all relevant levels of a color class as stated in \cref{def:generalized_separation}. As the arrangement has finitely many cells, we can place a large enough ball $B$ that contains all the bounded cells. Observe that we can assume $x_s$ to lie outside of~$B$. Consider a subset $I\subseteq [d]$ and a family of vectors $s^1,\ldots,s^k$ for which $s_i^1=\ldots =s_i^k$ for all $i\in I$. Let $\mathcal{F_{I}}$ be the union of all relevant levels, that is, the $\alpha_i$-levels for $\alpha_i \in \{\beta_i, \ldots, \gamma_i\}$, of the color classes $i \in I$. We claim that the points $x_{s^1},\ldots,x_{s^k}$ lie in a common unbounded cell of $\mathcal{F_{I}}$. Indeed, assume for the sake of contradiction that two of the points $x_{s^i}$ and $x_{s^j}$ lie in different unbounded cells. Then there is a level of $\mathcal{F_{I}}$ separating them outside of $B$. However, by definition, $x_{s^i}$ and $x_{s^j}$ must lie on the same side of any relevant level in $\mathcal{F_{I}}$.

We thus indeed have that all the points $x_{s^1},\ldots,x_{s^k}$ lie in a common unbounded cell of $\mathcal{F_{I}}$ outside of $B$. 
Now we use the same argument as in the proof of \cref{thm:rainbowHS}.
In particular, we can again find an embedding of the cube $[-1,1]^d$ where each face $F_I$ characterized by a subset $I \subseteq [d]$ 
lies in this common unbounded cell of $\mathcal{F_{I}}$. After applying a homeomorphism on $\mathbb{R}^d$ that maps this embedding to the cube $[-1,1]^d\subset\mathbb{R}^d$ we can thus assume that all relevant levels of the color class $\mathcal{F}_i$ have the facet $x_i=-1$ on their negative side, and the facet $x_i=1$ on their positive side.

We now take the same functions $f_i$ as in the proof of Theorem \ref{thm:rainbowHS}, which satisfy the conditions of the Poincar\'{e}-Miranda theorem.
We thus get that there is a point $x^*$ in $[-1,1]^d$ for which $f_i(x^*)=0$ for all $i\in\{1,\ldots,d\}$.
\end{proof}

\end{document}
